% ===========================================================================
% CAPÍTULO 2 — MARCO TEÓRICO: el transporte de datos entre el PS y la PL
% ===========================================================================

\section{Transporte de datos entre el PS y la PL} \label{sec:teoria_transporte}

Implementar los transceptores serie en la lógica programable resuelve la mitad del
problema: la serialización de los bits. Queda la otra mitad, que es la que realmente
determina el rendimiento del sistema completo: \textbf{por dónde y a qué coste viajan los
bytes entre la memoria del procesador y esos transceptores}. El MPSoC ofrece varios
caminos para ello, y la elección entre ellos no es un detalle de implementación, sino una
decisión de arquitectura con consecuencias directas sobre la carga de CPU, la latencia y
la ocupación de la FPGA.

Esta sección describe los mecanismos disponibles. Su aplicación concreta se desarrolla en la
sección~\ref{sec:transporte_desarrollo}.

\subsection{Bus AXI}

La comunicación entre el PS y la PL del Zynq UltraScale+ se realiza a través de puertos que
implementan el estándar AMBA AXI4 de ARM \cite{arm2011axi4}. El estándar define tres
perfiles:

\begin{itemize}
    \item \textbf{AXI4-Lite.} Perfil reducido, orientado a la lectura y escritura de
    \textit{registros} individuales. Cada transacción mueve una palabra y requiere que el
    procesador ejecute una instrucción de carga o almacenamiento. Este bus es comúnmente utilizado
    para configurar periféricos.

    \item \textbf{AXI4 \textit{full} (\textit{memory-mapped}).} Perfil completo, con soporte
    de ráfagas (\textit{bursts}) de hasta 256 transferencias por transacción. Es el bus que
    utilizan los maestros de la PL para acceder a la DDR del PS a través de los puertos
    de alto rendimiento (HP), sin pasar por la CPU.

    \item \textbf{AXI4-Stream.} Perfil sin direcciones: un flujo continuo de datos con
    señales de \textit{handshake} (\texttt{TVALID}/\texttt{TREADY}), marca de fin de paquete
    (\texttt{TLAST}) y, opcionalmente, un identificador de destino (\texttt{TDEST}). Es el
    perfil idóneo para conectar un periférico que produce o consume bytes en secuencia,
    como un UART.
\end{itemize}


\subsection{Acceso directo a memoria (DMA)}

Un controlador DMA es un maestro de bus que transfiere bloques de datos entre la memoria y un
periférico \textbf{sin intervención de la CPU}. El procesador se limita a programar la
transferencia (dirección de origen o destino y longitud) y a recibir una única interrupción
cuando el bloque completo ha terminado. La tasa de interrupción deja de depender del número
de bytes y pasa a depender del número de \textit{paquetes}, lo que en la práctica supone una
reducción de dos a tres órdenes de magnitud.

En el ecosistema de Xilinx, dos IP cubren este papel.

\subsubsection{AXI DMA (PG021)}

El IP \texttt{axi\_dma} \cite{xilinx2022pg021} traduce entre AXI4-Stream y AXI4
\textit{memory-mapped} en dos canales independientes:

\begin{itemize}
    \item \textbf{MM2S} (\textit{Memory-Mapped to Stream}): lee de la memoria y emite un flujo
    AXI-Stream hacia el periférico.
    \item \textbf{S2MM} (\textit{Stream to Memory-Mapped}): recibe un flujo del periférico y lo
    escribe en memoria. El fin de paquete lo marca la señal
    \texttt{TLAST} del \textit{stream}.
\end{itemize}

Cada canal admite dos modos de funcionamiento. En modo \textbf{directo} (\textit{simple DMA})
el procesador escribe la dirección del \textit{buffer} y la longitud en sendos registros, y el
motor ejecuta una sola transferencia. En modo \textbf{\textit{scatter-gather}} (SG) el
procesador construye en memoria una lista enlazada de \textit{descriptores de buffer} (BD),
cada uno con la dirección, la longitud y los bits de control del bloque, y el DMA la recorre
autónomamente. El modo SG permite encadenar transferencias sin que la CPU intervenga entre
ellas, a costa de que el propio DMA lea y escriba los descriptores en memoria.

La limitación de este IP es que \textbf{sirve a un único \textit{stream}}. Para $N$ periféricos
hacen falta $N$ instancias, con el consiguiente coste de área que esto supone, además de la
complejidad en software de gestionar $N$ controladores independientes y $N$ interrupciones. 

\subsubsection{AXI MCDMA (PG288)}

El IP \texttt{axi\_mcdma} \cite{xilinx2022pg288} es la respuesta a esa limitación: un
controlador \textbf{multicanal} que multiplexa hasta 16 canales lógicos sobre un único motor y
un único puerto hacia la memoria. Cada canal tiene su propio juego de registros, su propio
anillo de descriptores y su propio \textit{buffer} en la DDR, pero comparten el hardware de
transferencia.

El encaminamiento entre canales lógicos y flujos físicos se realiza gracias a la señal
\textbf{\texttt{TDEST}} del AXI-Stream. En el sentido S2MM (entrada), el MCDMA lee el
\texttt{TDEST} que acompaña a cada paquete, lo usa como identificador de canal y escribe la
carga útil en el \textit{buffer} apuntado por el descriptor activo de ese canal; el
demultiplexado, por tanto, lo resuelve el propio hardware sin intervención del procesador. En
el sentido MM2S (salida), la operación es la inversa: el motor etiqueta cada paquete que emite
con el \texttt{TDEST} del canal que lo originó, de modo que el periférico o el conmutador
\textit{stream} situado aguas abajo pueda encaminarlo a su destino. Basta así un único par de
interfaces AXI-Stream para dar servicio a los $N$ periféricos, frente a las $N$ instancias
independientes que exigiría el \texttt{axi\_dma}.


